\chapter{Syntax: The Modular Brick-Stack}

\section{The Rigid Subject-Object-Verb (SOV) Template}
PED maintained a non-negotiable SOV word order. This structural rigidity is what allowed for the preservation of the syntactic signal across 40,000 years. Because the 'Subject' was always an Active Agent and the 'Object' was always an Inactive Target, the relational geometry of the sentence remained fixed even as the individual sounds drifted.

\section{The Agglutinative Brick-Stacking Model}
Meaning in PED was not achieved through inflection (modifying the root), but through \textbf{Agglutinative Stacking}. A sentence was constructed like a wall of bricks, where each suffix added a discrete layer of logical information.

\subsection{The Universal Suffix Stack}
The ordering of the 'bricks' was deterministic:
\begin{center}
\textbf{[ROOT] + [ANIMACY] + [TARGET] + [PERSON]}
\end{center}

\section{Converbal Logic: The Connectivity of Thought}
PED lacked subordinating conjunctions such as "if," "when," or "because." Complex logical relationships were handled by applying noun-case markers directly to the verb stem, creating 'Converbs'.
\begin{itemize}
    \item \textbf{The Locative Converb (*-R):} "In/At [Action]" $\rightarrow$ "While/When [Action]."
    \item \textbf{The Dative Converb (*-K):} "Toward [Action]" $\rightarrow$ "In order to [Action]."
    \item \textbf{The Conditional Converb (*-i):} "Assuming [Action]" $\rightarrow$ "If [Action]."
\end{itemize}

\section{The Super-Sentence: Syntactic Reconstruction}
We provide a full syntactic gloss of a high-complexity PED sentence:
\begin{center}
\textit{*eK a-R T'w-N T'ar-N K'L-S-mi} \\
\small [1.SG.ACT DIST-LOC TWO-INACT TREE-INACT SEE-PFV-1.SG] \\
\normalsize "I saw those two trees."
\end{center}

\subsection{Reflex Derivation Table}
\begin{table}[h]
\centering
\begin{tabular}{lll} \toprule
\textbf{Language} & \textbf{Sentence Fragment} & \textbf{Structural Preservation} \\ \midrule
Vedic & \textit{aham tat dwe dāru paśyāmi} & PRON + DEM + NUM + SOV \\
Old Tamil & \textit{yāṉ atu iraṇṭu tāram kaṇṭēṉ} & PRON + DEM + NUM + SOV \\
PED Root & \textit{*eK a-R T'w-N T'ar-N K'LS-mi} & Pure Modular Stack \\ \bottomrule
\end{tabular}
\end{table}

\section{Sandhi: Phonetic Friction in the Stack}
As the 'bricks' were stacked together, the boundaries between suffixes often blurred due to phonetic friction (Sandhi). The most common rule in PED was the \textbf{Nasal Assimilation Rule}:
\begin{center}
\textit{*N + *K $\rightarrow$ *ŋ}
\end{center}
This explains why many internal clusters in Sanskrit and Tamil appear with a velar nasal before a dative or directional marker.
